%% =====================================================================
%% PWM-LDCT Dataset Paper — Nature Scientific Data submission
%%
%% Target venue: Nature Scientific Data (Data Descriptor format)
%% Template: Springer Nature; in development we use article + sn-style;
%%           switch to sn-jnl.cls (Springer Nature Journals template)
%%           at submission. The journal accepts LaTeX submissions through
%%           the Snapp / Latex.online system.
%%
%% Status: v0.1 working draft. RESULTS sections marked with \todo{} are
%%         placeholders to fill as Phase 1/2/3 of WS-1 lands data.
%%
%% Build:  pdflatex manuscript && bibtex manuscript && pdflatex manuscript
%%         (×2 for cross-refs)
%% =====================================================================

\documentclass[twoside,11pt]{article}

%% --- Packages ---------------------------------------------------------
\usepackage[utf8]{inputenc}
\usepackage[T1]{fontenc}
\usepackage{lmodern}
\usepackage{microtype}
\usepackage[margin=1in]{geometry}
\usepackage{graphicx}
\usepackage{booktabs}
\usepackage{multirow}
\usepackage{amsmath,amssymb}
\usepackage{siunitx}
\usepackage{xcolor}
\usepackage[colorlinks=true,linkcolor=blue,citecolor=blue,urlcolor=blue]{hyperref}
\usepackage{xurl}            % better URL line-breaking inside \url{}
\usepackage[round]{natbib}
\usepackage{authblk}
\usepackage{lineno}
\linenumbers

%% --- TODO macros ------------------------------------------------------
\newcommand{\todo}[1]{\textcolor{red}{\textbf{[TODO: #1]}}}
% Compile-clean placeholder floats. Replace each with the real table /
% figure as data lands; the \label{} keeps cross-references working.
\newcommand{\todoTable}[2]{%
  \begin{table}[ht]\centering%
    \fbox{\parbox{0.85\linewidth}{\centering\todo{TABLE: #2}}}%
    \caption{#2 (placeholder; populate when data lands).}\label{#1}%
  \end{table}}
\newcommand{\todoFigure}[2]{%
  \begin{figure}[ht]\centering%
    \fbox{\parbox{0.85\linewidth}{\centering\todo{FIGURE: #2}}}%
    \caption{#2 (placeholder; populate when data lands).}\label{#1}%
  \end{figure}}
\newcommand{\Nscans}{\todo{N\textsubscript{paired scans}}}
\newcommand{\Npatients}{\todo{N\textsubscript{patients}}}
\newcommand{\NvendorA}{\todo{N\textsubscript{vendor A}}}
\newcommand{\NvendorB}{\todo{N\textsubscript{vendor B}}}

%% --- Title block (Nature Scientific Data Data Descriptor) -------------
\title{\textbf{PWM-LDCT: A Multi-Vendor, Multi-Site, Paired-Dose CT Benchmark Dataset \\ for Reproducible Low-Dose Reconstruction Research}}

\author[1,2]{\textit{Authors to be confirmed at submission}}
\affil[1]{University of Texas Southwestern Medical Center, Dallas, TX, USA}
\affil[2]{PWM Protocol Foundation}

\date{Working draft v0.1 --- \today}

\begin{document}
\maketitle

%% =====================================================================
\begin{abstract}
\noindent We present \textbf{PWM-LDCT}, a multi-vendor, multi-site, paired-dose chest- and abdominal-CT dataset designed as a permanent, reproducibility-first benchmark for low-dose CT reconstruction research. The dataset comprises \Nscans{} paired full-dose / low-dose scans across at least two manufacturer platforms (\todo{vendor list}) and \todo{N$_{\textrm{sites}}$} clinical sites; \NvendorA{} patients contribute real reduced-tube-current paired acquisitions and \NvendorB{} contribute physically calibrated Poisson-thinned low-dose simulations alongside their full-dose scans. Each release record ships with raw sinograms, reconstructed images at four dose levels (10\%, 25\%, 50\%, 100\%), majority-vote radiologist annotations for lesion detection and segmentation tasks, and full provenance metadata under PhysioNet credentialed access. A pip-installable Python loader (\texttt{pwm\_ldct\_loader}) and three reproducible Docker pipelines deliver bit-identical preprocessing across machines, and the dataset's manifest is uniquely identified by a SHA-256 content hash that is mirrored to multiple immutability anchors (PhysioNet DOI, IPFS replica, and a PWM L3 specification registry entry) so that leaderboard submissions and downstream citations can be tied to a specific frozen benchmark version regardless of which anchor is consulted. By foregrounding multi-vendor diversity, real paired low-dose acquisitions, task-level annotations, and content-addressed version control, PWM-LDCT addresses three long-standing gaps in the field: vendor-specificity of published methods, evaluation through PSNR alone rather than clinical task performance, and benchmark drift across publication cycles. PWM-LDCT is designed as the empirical substrate for the recurring PWM Low-Dose CT Challenge and as a validation cohort for a companion signal-equivalence evaluation framework (both in preparation,~\citealt{ws4leaderboard,ws2framework}); the descriptor's claims (multi-vendor coverage, real paired-dose acquisitions, content-addressed manifest, annotation reliability, and reproducibility) stand independently of those companion artifacts.
\end{abstract}

%% =====================================================================
\section*{Background \& Summary}

Low-dose computed tomography (CT) reconstruction has been transformed by deep learning over the past decade~\citep{chen2017redcnn,wang2020deeplearningreview}, with successive methods reporting incremental gains in peak signal-to-noise ratio (PSNR) and structural similarity (SSIM) against published benchmarks~\citep{aapm2016challenge}. Yet three structural problems prevent the field from converging on a single, citable state-of-the-art:

\begin{enumerate}
    \item \textbf{Vendor specificity.} Most published methods are evaluated on data from one CT manufacturer, typically Siemens (via the AAPM 2016 Low-Dose CT Grand Challenge~\citep{aapm2016challenge}) or a vendor-internal dataset. A method that wins on Siemens data may degrade by 2--4~dB PSNR on GE or Canon data because of subtle differences in raw-projection geometry, vendor-specific reconstruction kernels, and noise statistics. Without a multi-vendor test split, cross-vendor robustness cannot be measured.
    \item \textbf{PSNR / SSIM are not clinical endpoints.} A reconstruction that scores well on pixel-wise metrics may smooth out clinically relevant low-contrast lesions~\citep{kim2019smoothing}. Without paired task annotations (lesion bounding boxes, segmentations, diagnostic-quality labels) reviewers cannot ask the diagnostically important question: does the low-dose reconstruction still support correct clinical decisions?
    \item \textbf{Benchmark drift.} As datasets evolve (new patients added, splits reshuffled, preprocessing changed), reported numbers from year~$n$ become incomparable to year~$n+1$. The community lacks a frozen, content-addressed reference version to anchor longitudinal comparison.
\end{enumerate}

PWM-LDCT is designed to close all three gaps simultaneously. Multi-vendor coverage is enforced by inclusion criteria; the corresponding subgroup statistical power is justified explicitly (Section \textit{Sample size justification}). Per-patient majority-vote radiologist annotations support a curated set of clinical tasks (lung nodule detection, liver lesion segmentation, abdominal organ quantification) collected under a documented multi-reader QC protocol (Section \textit{Annotation protocol}). The dataset's manifest is content-addressed by a SHA-256 hash that downstream users can verify offline; the hash is mirrored to PhysioNet (DOI), IPFS, and the PWM L3 specification registry (\texttt{pwm\_l3\_spec:pwm\_low\_dose\_ct\_benchmark:1.0.0}), so any future paper citing the hash refers to bit-identical content regardless of which anchor remains available. An explicit versioning policy (Section \textit{Data Versioning and Persistent Identifiers}) governs future extensions and corrections.

Table~\ref{tab:dataset_comparison} situates PWM-LDCT against existing public datasets. LIDC-IDRI~\citep{lidc_idri} is the field's largest annotated chest cohort but is single-dose and single-vendor in practice. The AAPM 2016 Low-Dose CT Grand Challenge release~\citep{aapm2016challenge} is the only public source of real paired-dose acquisitions but is limited to 10 patients on a single Siemens platform. The Mayo Clinic LDCT-and-Projection-Data collection~\citep{mccollough2020mayo} extends AAPM 2016 with raw projections but inherits its single-vendor scope. PWM-LDCT is the first multi-vendor paired-dose release at the $\geq$ 500-patient scale and the first with a content-addressed benchmark hash anchored to a public registry.

\begin{table}[ht]
    \centering
    \small
    \setlength{\tabcolsep}{4pt}
    \begin{tabular}{lccccc}
        \toprule
        \textbf{Dataset} & \textbf{Patients} & \textbf{Vendors} & \textbf{Real paired LD} & \textbf{Projections} & \textbf{Multi-anchor content hash} \\
        \midrule
        LIDC-IDRI~\citep{lidc_idri}       & 1{,}018 & 1 (mostly)    & No  & No   & No \\
        AAPM 2016~\citep{aapm2016challenge} & 10      & 1 (Siemens)   & Yes & Yes  & No \\
        Mayo LDCT-PD~\citep{mccollough2020mayo} & 299  & 1 (Siemens)   & Yes & Yes  & No \\
        \textbf{PWM-LDCT (ours)}              & \Npatients & $\geq$ 2 & Yes & Yes  & \textbf{Yes (DOI / IPFS / L3)} \\
        \bottomrule
    \end{tabular}
    \caption{Comparison of PWM-LDCT against existing public low-dose CT datasets. ``Real paired LD'' indicates whether the dataset contains real (not simulated) paired full-dose / low-dose acquisitions for the same patient. ``Multi-anchor content hash'' indicates whether the dataset version is identified by a SHA-256 manifest hash mirrored across multiple independent persistent-identifier anchors so that long-term citation resolves consistently.}
    \label{tab:dataset_comparison}
\end{table}

The dataset is open to academic and industrial researchers under PhysioNet credentialed access (HIPAA Safe Harbor compliant for de-identified scans). Companion methodological work --- a reference reconstruction method~\citep{ws3reference}, a probabilistic signal-equivalence evaluation framework~\citep{ws2framework}, and a permanent leaderboard~\citep{ws4leaderboard} --- is in preparation and signposted from Usage Notes; this descriptor's scientific claims (multi-vendor coverage, real paired-dose acquisitions, content-addressed manifest, annotation reliability, and dataset reproducibility) stand independently of those companion artifacts.

%% =====================================================================
\section*{Methods}

The data collection, processing, annotation, and release workflow is summarized in Figure~\ref{fig:workflow}.

\todoFigure{fig:workflow}{Workflow diagram. Rendered figure should show six sequential stages connected by arrows: (1) patient enrollment and consent; (2) paired full-dose / real-low-dose acquisition (or full-dose-only for the simulation-augmented cohort); (3) two-stage de-identification (DICOM tag stripping + Tesseract OCR sweep); (4) reconstruction at four dose levels (real or Poisson-thinned simulation, plus FBP / vendor IR / TV variants); (5) multi-reader annotation with majority-vote ground truth + per-scan Likert diagnostic-quality scoring; (6) quality control + SHA-256 manifest + PhysioNet deposit + content-hash mirroring to PWM L3 registry and IPFS. Side annotations should call out the IRB approval, the consent waiver branch for retrospective cohorts, the cross-vendor harmonization step between stages 2 and 4, and the inter-rater audit feedback loop between stages 5 and 6.}

\subsection*{Ethics, consent, and data sharing}

The study was reviewed and approved by the University of Texas Southwestern Medical Center Institutional Review Board (protocol \todo{IRB \#}), under which UTSW served as the coordinating site for all participating institutions. Partner sites operated under their respective institutional IRB approvals plus a UTSW-led data sharing agreement (DSA) governing transfer of de-identified data to the coordinating site for release preparation.

\paragraph{Consent.} Two consent paths were used. \textbf{Prospectively enrolled patients} (the cohort contributing real paired-dose acquisitions) provided written informed consent for the use of their de-identified imaging data, including raw projection data and reconstructed images, in research and for redistribution under credentialed access. \textbf{Retrospective patients} (the cohort contributing full-dose scans used as the substrate for low-dose simulation) were included under an IRB-approved waiver of consent (45~CFR~46.116(f)) on the grounds that (a) the research presents no greater than minimal risk, (b) the waiver does not adversely affect the rights and welfare of the subjects, (c) the research could not practicably be carried out without the waiver, and (d) subjects will be provided with no individually-attributable information.

\paragraph{Data sharing committee.} A standing data sharing committee comprising representatives from each contributing site reviews all credentialed-access requests received via PhysioNet. The committee operates on a default-approve policy for academic and industrial researchers who agree to the standard data use agreement (DUA); requests proposing redistribution, re-identification attempts, or commercial-product training without a separate commercial license are referred to UTSW Office of Technology Development for individual review.

\paragraph{Subject withdrawal.} Patients in the prospective cohort retain the right to withdraw consent at any time; withdrawn patients' scans are removed from the next dataset version with a tombstone entry retained in the version-control record for accountability. Because the dataset is content-addressed and immutable per version (Section \textit{Data Versioning and Persistent Identifiers}), withdrawal triggers a minor-version increment with the withdrawn patient excluded from the new version's manifest. Prior-version downloads cannot be retroactively edited but the canonical pointer is advanced.

\paragraph{Conflict-of-interest disclosure.} The PWM Protocol Foundation provided no direct funding to the data-collection sites; foundation-affiliated authors hold no financial interest in any specific reconstruction method submitted to the leaderboard. The complete COI statement is in the Competing Interests section.

\subsection*{Cohort recruitment and acquisition protocol}

PWM-LDCT scans were acquired prospectively at \todo{N$_{\textrm{sites}}$} clinical sites (University of Texas Southwestern Medical Center serving as the coordinating site) between \todo{collection start date} and \todo{collection end date}. Inclusion criteria were:

\begin{itemize}
    \item Adult patients (age $\geq$ 18) undergoing routine chest CT for lung-cancer screening or routine abdominal CT for oncology follow-up
    \item Informed consent obtained for the use of de-identified imaging data in research
    \item Acquisition on a manufacturer-platform combination not represented at $\geq$ 30\% within the cohort (enforcing vendor diversity)
\end{itemize}

For each enrolled patient a \textit{paired} acquisition was performed:

\begin{itemize}
    \item \textbf{Full-dose scan}: the patient's clinically indicated tube current setting (typically 200--400~mA depending on body habitus).
    \item \textbf{Real low-dose scan}: an immediately subsequent acquisition at 25\% of the full-dose mA, taken within the same breath-hold for chest scans and within $\leq$ 30~seconds for abdominal scans to minimize motion-induced misregistration. This paired acquisition was performed on \NvendorA{} patients across vendors that support rapid sequential acquisition at differing mA.
    \item \textbf{Simulated low-dose data}: for \NvendorB{} additional patients (and as augmentation for the real-paired cohort) we generated 10\%, 25\%, and 50\% dose-equivalent sinograms via physically calibrated Poisson-noise insertion on the full-dose sinogram, using the forward model implemented in PWM's \texttt{ct\_radon} modality library~\citep{pwmcore}. The simulation protocol is documented in Section \textit{Low-dose simulation} below; we explicitly maintain both real and simulated low-dose images in the dataset so that downstream users can assess how well methods trained on simulated low-dose generalize to real low-dose.
\end{itemize}

\subsection*{Radiation dose justification for the paired-acquisition cohort}

The \NvendorA{} prospectively-enrolled patients in the paired-acquisition cohort received an additional research-purpose low-dose acquisition beyond the clinically-indicated full-dose scan. The IRB review, consent path, and data-sharing governance for this cohort are documented in Ethics, consent, and data sharing above; this subsection adds the radiation-dose-specific justification that reviewers and oversight bodies expect for any added-dose research protocol.

\paragraph{Scientific necessity.} The principal scientific aim of the paired-acquisition cohort is to ground-truth low-dose CT simulation against physical low-dose acquisition at multi-vendor scale. Simulated low-dose imagery, generated by Poisson thinning of full-dose sinograms, is the substrate of most published low-dose CT reconstruction research, yet the systematic discrepancy between simulation and physical reduced-tube-current acquisition has not been characterized at a sample size sufficient to be statistically informative across vendors~\citep{wang2020deeplearningreview}. Without paired physical-acquisition data, simulation-trained methods cannot be validated against real low-dose noise, and the field cannot answer whether simulation-bench-leading methods transfer to clinical use. We considered but rejected a purely retrospective design because (i) clinically-indicated mA reduction is confounded with body-habitus-driven mA modulation that frustrates the simulation comparison, and (ii) no existing public release pairs full-dose and physical low-dose acquisitions on the same patient at the $\geq$ 500-patient scale across vendors. The McCollough et al.\ paired-dose acquisition protocol on Siemens scanners~\citep{mccollough2020mayo} is the closest methodological precedent and is the template against which our multi-vendor cohort is designed.

\paragraph{Additional dose budget.} The additional 25\%-of-full-dose mA acquisition delivers an estimated \todo{$\Delta\textrm{CTDI}_{\textrm{vol}}$ (mGy, chest)} additional dose per patient for chest studies and \todo{$\Delta\textrm{CTDI}_{\textrm{vol}}$ (mGy, abdomen)} for abdominal studies, corresponding to an estimated effective-dose increment of \todo{$\Delta E$ (mSv, chest)} and \todo{$\Delta E$ (mSv, abdomen)} respectively. For context, this increment is approximately \todo{X}\,\% of the typical clinically-indicated full-dose study and approximately \todo{Y}\,\% of typical annual background radiation ($\sim$~3\,mSv in the US). Per-patient dose values are recorded directly from the scanner's DICOM RDSR (Radiation Dose Structured Report) and surfaced in \texttt{metadata.json}, so downstream users can verify the dose budget per patient rather than rely on protocol-level estimates.

\paragraph{ALARA compliance.} The 25\%-of-full-dose mA setting was selected as the most informative paired dose for benchmark utility while remaining below the additional-dose ceiling negotiated with the coordinating IRB. No patient received more than one paired research-purpose low-dose acquisition. Patients with a cumulative recent imaging dose history above an institutional ceiling (\todo{cumulative dose ceiling}) were excluded, as were pregnant patients and patients $<$ 18 years of age. Tube-current modulation was held constant between the full-dose and paired acquisitions to keep the dose-reduction comparison clean; this is documented in the per-scan acquisition log preserved in \texttt{metadata.json}.

\paragraph{Dose disclosure in the informed-consent script.} Consent for the additional research-purpose acquisition was obtained as part of the written informed-consent process described in Ethics, consent, and data sharing. The IRB-approved consent script explicitly disclosed (i) the research purpose of the second acquisition, (ii) the estimated additional radiation dose (numerical estimates above), (iii) the absence of direct diagnostic benefit from the additional acquisition, and (iv) the patient's right to refuse the additional acquisition with no effect on their clinical care. The consent template is published as \texttt{schema/research\_consent\_template.md} in the project repository.

\subsection*{Sample size justification}

The dataset's target of $\geq$ 500 paired patients was set by a standard statistical-power analysis for the canonical clinical task (lung-nodule detection AUC at fixed false-positive rate). For a paired-bootstrap two-method comparison with target equivalence margin $\varepsilon$ at significance $\alpha$, the leading-order sample-size requirement is

\begin{equation}
    n \geq \left(\frac{z_{1-\alpha/2}}{\varepsilon}\right)^2 \sigma^2 (1 - \rho)
\end{equation}

where $\sigma$ is the per-patient performance-difference standard deviation and $\rho$ is the pairing correlation between the two arms. For the canonical target $(\varepsilon = 0.02,\ \alpha = 0.05)$ on AUC with conservative $\sigma \approx 0.10$ and $\rho \approx 0.5$, this yields $n \geq 192$ patients per evaluation slice. Because the dataset must support \textit{multiple} stratified evaluations simultaneously --- per dose level, per vendor, per anatomy, per task --- the realized cohort target was scaled to $\geq 500$ patients to retain statistical power at each subgroup. The corresponding cohort-stratification budget is summarized in Table~\ref{tab:power_budget} \todo{populate at submission}. Methods that wish to additionally compute 5-tuple credentials under the companion signal-equivalence framework (manuscript in preparation,~\citealt{ws2framework}) will find that the cohort meets that framework's matching leading-order requirements at canonical target $(\varepsilon = 0.02,\ \alpha = 0.05)$, but the descriptor's sizing argument is framework-independent.

\todoTable{tab:power_budget}{Sample-size budget by stratification: rows are evaluation conditions (e.g., 25\%-dose chest, 25\%-dose abdomen, leave-one-vendor-out cross-vendor); columns give required $n$ at $(\varepsilon, \alpha) = (0.02, 0.05)$ vs realized $n$ in the released cohort}

\subsection*{Low-dose simulation}

We simulate low-dose acquisitions from full-dose sinograms using the physically-calibrated forward model implemented in the open-source \texttt{pwm\_core.contrib.modalities.ct\_radon} package, following standard practice in the low-dose CT literature~\citep{zeng2015nonlocal,wang2020deeplearningreview}. The model combines Poisson photon-counting noise with an additive electronic-noise floor and ray-dependent incident-photon-count weighting:

\begin{equation}
    N_r(\ell) \sim \textrm{Poisson}\!\big(I_0(\ell) \cdot r \cdot e^{-s_{\textrm{ref}}(\ell)}\big) + \mathcal{N}\!\big(0,\sigma_e^2\big), \qquad \tilde{s}_{\textrm{low}}(\ell) = -\ln\!\left( \frac{\max\!\big(N_r(\ell),\,1\big)}{I_0(\ell) \cdot r} \right)
    \label{eq:lowdose_sim_v1}
\end{equation}

where $s_{\textrm{ref}}(\ell)$ is the full-dose effective-monochromatic log-attenuation projection along ray $\ell$; $I_0(\ell) = I_0^{(0)} \cdot b(\ell)$ is the ray-dependent incident-photon count combining a per-scanner monochromatic-equivalent reference $I_0^{(0)}$ (calibrated from the acquisition's recorded tube parameters and manufacturer-published kVp / mAs / pitch tables) with the scanner's bowtie-filter profile $b(\ell)$ (calibrated from manufacturer-published or in-house phantom measurements); $\sigma_e$ is the per-scanner electronic-noise standard deviation, calibrated from low-flux phantom measurements~\citep{zeng2015nonlocal}; and $r \in \{0.10, 0.25, 0.50\}$ is the dose-reduction ratio. The $\max(\cdot,1)$ guard prevents log-of-zero at extreme attenuations.

\paragraph{Simulation approximations and known limitations.} Equation~\eqref{eq:lowdose_sim_v1} is an effective-monochromatic forward model and does not capture all physical effects of polychromatic CT acquisition. We make four approximations explicit so users can judge applicability:

\begin{itemize}
    \item \textbf{Monochromatic approximation.} We use a kVp-equivalent effective mean photon energy rather than integrating over a polychromatic X-ray spectrum. This is the standard simplification in published LDCT simulators~\citep{zeng2015nonlocal,wang2020deeplearningreview}; it does not capture beam-hardening artifacts, which therefore appear with the same intensity in both simulated and full-dose images and are not differentially induced by dose reduction.
    \item \textbf{Bowtie + spectrum from manufacturer tables.} Where DICOM headers do not preserve the per-scan bowtie profile, we fall back to a uniform $b(\ell) \equiv 1$ and flag the affected scans in \texttt{metadata.json}; users can opt to drop these scans from a sensitivity analysis.
    \item \textbf{Detector quantum efficiency.} Detector QE is folded into the calibrated $I_0^{(0)}$ rather than modeled separately. Users requiring explicit DQE modeling (e.g., for detector-comparison studies) should not use these simulated images.
    \item \textbf{Electronic-noise floor.} The Gaussian electronic-noise term $\mathcal{N}(0,\sigma_e^2)$ is calibrated globally per scanner rather than spatially-varying across the detector array. Position-dependent electronic noise (a $\leq$~5\,\% effect on typical clinical scanners) is not modeled.
\end{itemize}

The discrepancy between simulated and real low-dose at $r = 0.25$ is quantified per-vendor on the real-paired-dose cohort in the Technical Validation section (Figure~\ref{fig:sim_vs_real}); users running simulation-vs-real sensitivity studies should consult that figure. The forward operator and Poisson sampling are implemented as the open-source \texttt{pwm\_core.contrib.modalities.ct\_radon} package; we explicitly do not reinvent the forward model so any downstream user replicating the simulation pipeline obtains bit-identical sinograms.

\subsection*{Reconstruction}

For each (patient, dose-level) pair we provide three image-space reconstructions:

\begin{enumerate}
    \item \textbf{Filtered back-projection (FBP)} with a Hamming-windowed Ram-Lak filter --- the canonical analytical reference.
    \item \textbf{Vendor standard iterative reconstruction (IR)} --- the algorithm clinically deployed on the originating scanner, where the vendor exposes a research interface (Siemens \textit{Admire}, GE \textit{Veo}, Canon \textit{AIDR3D}). Where unavailable, the field is recorded as \texttt{null}.
    \item \textbf{Total-variation (TV) regularized reconstruction} --- the open reference iterative method used to ensure that even vendor-restricted institutions have a reproducible iterative baseline.
\end{enumerate}

All reconstructions are provided alongside the raw projections so that researchers can apply their own reconstruction algorithms to the same sinograms.

\subsection*{Cross-vendor harmonization}

A multi-vendor benchmark is only useful if cross-vendor comparison is methodologically defensible. We applied four explicit harmonization steps so that scans from Siemens, GE, and Canon platforms are directly comparable while preserving vendor-specific characteristics that downstream methods may wish to model.

\paragraph{Hounsfield-unit calibration.} All vendors report HU on the standard scale, but small phantom-derived calibration offsets exist between platforms. We re-verified HU calibration at each site using the AAPM CT performance phantom and report the per-vendor offset (mean and 95\% CI) in the metadata so users can apply or skip a calibration correction as their method requires.

\paragraph{Reconstruction-kernel matching.} Vendors use platform-specific reconstruction kernel names (Siemens \textit{B30f}, GE \textit{Standard}, Canon \textit{FC18}) with similar but not identical frequency response. Per-vendor measured modulation transfer functions (MTFs) at $\textrm{MTF}_{50}$ and $\textrm{MTF}_{10}$ are recorded in \texttt{metadata.json} for each scan to make the kernel choice visible to downstream methods. We deliberately did \emph{not} re-reconstruct all scans to a common kernel, because doing so would discard the very vendor-specific structure that makes a multi-vendor dataset useful.

\paragraph{Slice-thickness handling.} Slice thickness varies clinically by indication (1.0--3.0~mm). The released DICOM preserves the acquired thickness; the loader (\texttt{pwm\_ldct\_loader}) exposes a \texttt{target\_slice\_thickness} parameter that interpolates to a common thickness on the fly. Downstream methods that depend on a fixed thickness can opt in; methods that explicitly model varying thickness can opt out.

\paragraph{Pixel-spacing normalization.} In-plane pixel spacing varies across scanners and reconstruction settings (0.5--1.0~mm). Released images preserve original pixel spacing; the loader provides an optional bilinear resampling to a common pixel grid (default 0.75~mm). The original spacing is preserved in the metadata so that physics-respecting reconstruction methods can recover the exact acquisition geometry.

\paragraph{What harmonization is \emph{not} applied.} We do not normalize noise statistics across vendors (different noise distributions are a feature, not a bug, for cross-vendor generalization studies). We do not normalize motion or breath-hold artifacts (these are clinical realities that downstream methods must handle). We do not exclude scans that fall outside any single vendor's typical parameter range (which would systematically remove edge cases that real deployment encounters).

\subsection*{Annotation protocol}

Annotations were collected from a panel of \todo{N} board-certified radiologists (\todo{N} fellowship-trained in thoracic imaging, \todo{N} in abdominal imaging) with $\geq$~5 years of post-training experience, distributed across the participating sites. The protocol covered three annotation classes plus a per-scan diagnostic-quality assessment.

\paragraph{Annotator training and calibration.} Before annotating release scans, every panel member completed a calibration session on a 20-case training set drawn from outside the release cohort. Each radiologist's calibration-set agreement with a senior reference (intraclass correlation for continuous attributes; Cohen's $\kappa$ for categorical) was required to exceed pre-specified thresholds (\todo{ICC threshold}, \todo{$\kappa$ threshold}) before any of their annotations were admitted to the release set. Failed calibration triggered a second session; persistent failure resulted in exclusion from the panel (\todo{N} excluded).

\paragraph{Blinding.} Each radiologist was blinded to (a) the dose level of the scan they were annotating --- the same patient's 10\%, 25\%, 50\%, and 100\% dose reconstructions were presented in random order across non-adjacent sessions --- and (b) any annotation produced by other radiologists on the same case until adjudication. Blinding to dose level was enforced by stripping or randomizing all dose-identifying DICOM tags in the annotation interface; blinding to other annotators was enforced by case-randomized assignment within the annotation queue.

\paragraph{Annotation classes.}

\begin{itemize}
    \item \textbf{Lung nodule detection (chest scans).} Bounding boxes for all pulmonary nodules $\geq$~3~mm in maximum diameter, annotated independently by $\geq$~2 panel radiologists. Per-nodule attributes recorded: maximum diameter (mm), location (lobe + segment), texture (solid / part-solid / non-solid), Lung-RADS-aligned category.
    \item \textbf{Liver lesion segmentation (abdominal scans).} Per-voxel segmentation masks for hepatic lesions $\geq$~5~mm in maximum diameter, with the same independent multi-reader protocol; per-lesion size and LI-RADS category recorded.
    \item \textbf{Diagnostic-quality label (per scan).} Each radiologist rated each reconstruction on a 5-point Likert scale for \textit{diagnostic confidence} ($1$ = nondiagnostic, $5$ = fully diagnostic), with per-dose-level scoring so downstream analyses can correlate computational reconstruction quality with radiologist-perceived diagnostic utility.
\end{itemize}

\paragraph{Dispute resolution.} Discordant per-case annotations (e.g., one reader detected a nodule that another did not, or readers disagreed on a categorical attribute) were referred to a third radiologist for adjudication. The third-reader vote determined the released ground-truth label; the original discordant annotations are preserved in the release as \texttt{annotations/<case>/raw\_reader\_\textit{n}.json} for downstream studies of annotation uncertainty.

\paragraph{Periodic agreement audit.} Inter-rater agreement was monitored monthly throughout the annotation campaign. Pre-specified action thresholds: if monthly Cohen's $\kappa$ for nodule detection fell below \todo{$\kappa$ floor} or Dice for liver-lesion segmentation fell below \todo{Dice floor}, the annotation queue was paused and a re-calibration session held with all active annotators before resuming. The audit log is published as supplementary material so downstream users can see the time-course of inter-rater agreement.

\paragraph{Adjudication-disagreement-flagged subset.} Cases where the third reader's adjudication disagreed with \textit{both} initial readers are retained in the release but flagged with \texttt{adjudication\_disagreement: true}. Downstream studies can choose to include or exclude this subset; we recommend including it for robustness studies and excluding it for high-confidence task evaluation. The flagged-case fraction per dose level is reported in the Technical Validation section.

\subsection*{De-identification and HIPAA compliance}

All scans were de-identified per HIPAA Safe Harbor (45~CFR~164.514(b)(2)). The DICOM tag whitelist and the burned-in-text scrubbing recipe are published as \texttt{schema/dicom\_cleaning\_spec.md} in the project repository. De-identification was verified by a two-stage pipeline: (i) automated PyDICOM tag stripping against the whitelist; (ii) OCR over every reconstructed slice using Tesseract~\citep{tesseract2007}, with any positive hit (e.g., date, name, MRN) flagged for manual review and reprocessing. Every scan in the released set passed both stages with no PHI detected.

\subsection*{DICOM conformance}

Released imaging data preserves each scan's original DICOM SOP-class identifiers as recorded by the scanner. Reconstructed image series are released under SOP class \texttt{CT Image Storage} (UID \texttt{1.2.840.10008.5.1.4.1.1.2}); projection-domain sinograms are released under each vendor's research-purpose \texttt{Raw Data Storage} SOP class (Siemens, GE, and Canon vendor-equivalent SOP classes are documented per-source in \texttt{schema/dicom\_conformance.md}). The released DICOM is fully conformant with the relevant Part-10 file format. Patient orientation (\texttt{Patient Position}, \texttt{Image Orientation Patient}) is preserved as acquired; we do not reorient. The de-identification whitelist preserves all geometry-relevant tags (pixel spacing, slice thickness and position, gantry tilt, reconstruction kernel name, table position) needed to reproduce a scan's physical geometry; the side-by-side mapping between retained DICOM tags and the loader-exposed \texttt{metadata.json} fields is published at \texttt{schema/dicom\_to\_metadata\_mapping.md} so downstream users can recover original coordinates and acquisition parameters for any released scan without re-downloading the source.

\subsection*{Patient-level data splits}

We partition patients (not scans) into 60\% train / 20\% validation / 20\% test, stratified by anatomical site, vendor, and dose-pair availability. A patient's full-dose, real low-dose, and simulated low-dose scans all reside in the same split. Split assignments are deterministic given the published random seed (\texttt{42}) and a patient's de-identified ID. The validation split is provided for hyperparameter tuning; the test split is reserved for final-evaluation reporting in published work.

%% =====================================================================
\section*{Data Records}

The dataset is distributed via PhysioNet at \todo{PhysioNet URL} under credentialed access. The full release contains:

\begin{table}[h]
    \centering
    \begin{tabular}{lp{8.5cm}}
        \toprule
        \textbf{Path} & \textbf{Contents} \\
        \midrule
        \texttt{patients/<patient\_id>/} & Per-patient directory \\
        \quad \texttt{full\_dose/} & DICOM, reconstructed image, sinogram \\
        \quad \texttt{low\_dose\_real/} & DICOM, image, sinogram (where available) \\
        \quad \texttt{low\_dose\_sim\_r010/} & Simulated 10\% dose \\
        \quad \texttt{low\_dose\_sim\_r025/} & Simulated 25\% dose \\
        \quad \texttt{low\_dose\_sim\_r050/} & Simulated 50\% dose \\
        \quad \texttt{annotations/} & JSON annotations (bbox, seg, Likert) \\
        \quad \texttt{metadata.json} & Vendor, model, acquisition parameters \\
        \texttt{splits/train.txt}, \texttt{val.txt}, \texttt{test.txt} & Patient ID lists \\
        \texttt{schema/dataset\_schema.md} & Canonical record schema \\
        \texttt{schema/dicom\_cleaning\_spec.md} & De-identification protocol \\
        \texttt{checksums.sha256} & SHA-256 hash of every file \\
        \texttt{LICENSE.txt} & Release license (PhysioNet credentialed) \\
        \bottomrule
    \end{tabular}
    \caption{Data record layout. Per-patient directories contain all paired acquisitions and the full per-patient annotation set; top-level files enumerate splits, schema, and integrity hashes.}
    \label{tab:data_records}
\end{table}

Aggregate counts and split structure are summarized in Table~\ref{tab:cohort}; demographic characterization is in Table~\ref{tab:demographics}; per-vendor acquisition parameters are in Table~\ref{tab:acquisition_params}.

\begin{table}[h]
    \centering
    \small
    \begin{tabular}{lcccc}
        \toprule
        \textbf{Split} & \textbf{Patients} & \textbf{Vendors} & \textbf{Sites} & \textbf{Anatomy} \\
        \midrule
        Training        & \todo{N} & \todo{N} & \todo{N} & chest, abdomen \\
        Validation      & \todo{N} & \todo{N} & \todo{N} & chest, abdomen \\
        Test            & \todo{N} & \todo{N} & \todo{N} & chest, abdomen \\
        \textbf{Total}  & \todo{N} & \todo{N} & \todo{N} & --- \\
        \bottomrule
    \end{tabular}
    \caption{Cohort composition by split. Final patient counts are filled at submission.}
    \label{tab:cohort}
\end{table}

\begin{table}[h]
    \centering
    \small
    \setlength{\tabcolsep}{4pt}
    \begin{tabular}{lccc}
        \toprule
        \textbf{Characteristic} & \textbf{Chest cohort} & \textbf{Abdomen cohort} & \textbf{Combined} \\
        \midrule
        Patients ($N$)                            & \todo{} & \todo{} & \todo{} \\
        Age, median (IQR), years                   & \todo{} & \todo{} & \todo{} \\
        Age range, years                           & \todo{} & \todo{} & \todo{} \\
        Sex, female (\%)                           & \todo{} & \todo{} & \todo{} \\
        BMI, median (IQR), kg/m$^2$                & \todo{} & \todo{} & \todo{} \\
        Race / ethnicity (\% White / Black / Hispanic / Asian / Other) & \todo{} & \todo{} & \todo{} \\
        Smoking history (\% current / former / never) & \todo{} & --- & \todo{} \\
        Clinical indication breakdown              & \todo{} & \todo{} & \todo{} \\
        Real-paired-dose patients (\%)             & \todo{} & \todo{} & \todo{} \\
        Simulation-only patients (\%)              & \todo{} & \todo{} & \todo{} \\
        \bottomrule
    \end{tabular}
    \caption{Demographic and clinical characterization of the released cohort. Race / ethnicity follow OMB-1997 categories as recorded at the originating site; ``Other'' includes American Indian / Alaska Native, Native Hawaiian / Pacific Islander, and ``unknown / declined.''}
    \label{tab:demographics}
\end{table}

\begin{table}[h]
    \centering
    \small
    \setlength{\tabcolsep}{4pt}
    \begin{tabular}{lccc}
        \toprule
        \textbf{Parameter} & \textbf{Siemens} & \textbf{GE} & \textbf{Canon} \\
        \midrule
        Scanner model(s)                          & \todo{} & \todo{} & \todo{} \\
        Tube voltage (kVp), median (range)        & \todo{} & \todo{} & \todo{} \\
        Tube current (mA), full-dose, median (IQR) & \todo{} & \todo{} & \todo{} \\
        Tube current (mA), low-dose, median (IQR)  & \todo{} & \todo{} & \todo{} \\
        Effective mAs (full-dose), median (IQR)    & \todo{} & \todo{} & \todo{} \\
        Pitch, median (range)                      & \todo{} & \todo{} & \todo{} \\
        Slice thickness (mm), median (range)       & \todo{} & \todo{} & \todo{} \\
        Reconstruction kernel(s)                   & \todo{e.g.\ B30f, B50f} & \todo{e.g.\ Standard, Lung} & \todo{e.g.\ FC18, FC51} \\
        CTDI$_{\textrm{vol}}$ full-dose (mGy), median & \todo{} & \todo{} & \todo{} \\
        CTDI$_{\textrm{vol}}$ low-dose (mGy), median  & \todo{} & \todo{} & \todo{} \\
        Patients ($N$)                             & \todo{} & \todo{} & \todo{} \\
        \bottomrule
    \end{tabular}
    \caption{Per-vendor acquisition parameter summary. CTDI$_{\textrm{vol}}$ is the volume CT dose index reported by the scanner. Per-scan exact values are in each scan's \texttt{metadata.json}.}
    \label{tab:acquisition_params}
\end{table}

%% =====================================================================
\section*{Technical Validation}

We performed four classes of technical validation:

\subsection*{Reconstruction sanity}
For every full-dose scan we recomputed FBP from the released sinogram and verified that the recomputed image matches the released reconstruction within $\leq$~\todo{FP tolerance}~HU at every voxel. This protects against an off-by-one indexing error between sinogram release and image release, which is a common silent failure mode of CT-dataset releases.

\subsection*{Real- vs simulated-low-dose comparison}
On the \NvendorA{} patients with paired real-low-dose acquisitions, we measured the per-patient \textit{distributional} agreement between the simulated 25\%-dose image (generated from the patient's own full-dose sinogram) and the real 25\%-dose acquisition. The agreement metrics are reported in Figure~\ref{fig:sim_vs_real}. The purpose is not to assert simulation equivalence (it is well known that simulated low-dose noise underestimates real low-dose noise in some regimes~\citep{wang2020deeplearningreview}) but to make the discrepancy visible so users can decide which split is appropriate for their evaluation.

\todoFigure{fig:sim_vs_real}{Per-patient distributional agreement between simulated 25\%-dose and real 25\%-dose acquisitions on paired-acquisition patients}

\subsection*{Inter-rater annotation reliability}
Cohen's $\kappa$ for lung nodule detection (per-radiologist, per-scan) and Dice for liver lesion segmentation are summarized in Table~\ref{tab:irr}. We report agreement broken down by dose level so that any tendency of low-dose images to depress annotator agreement is visible.

\todoTable{tab:irr}{Inter-rater annotation reliability ($\kappa$ for nodule detection; Dice for liver lesion segmentation), broken down by dose level}

\subsection*{Baseline reproductions}
To demonstrate that the dataset supports end-to-end training and evaluation of representative deep reconstruction methods through the provided loader and Docker pipelines without code modification, we reproduce four published baselines on the PWM-LDCT training split and report their test-split performance in Table~\ref{tab:baselines}. The baselines span method classes: RED-CNN~\citep{chen2017redcnn} (encoder-decoder CNN), a published transformer-based reconstruction method (\todo{select published method, e.g., CTformer or TransCT}), a published diffusion-based reconstruction method (\todo{select published method}), and a published unrolled-iterative reconstruction method (\todo{select published method, e.g., LEARN}). This section is intentionally not a benchmark ranking: relative method comparison is out of scope for a Data Descriptor and is the subject of separate companion work. The role of these reproductions is solely to provide evidence that (i) the loader and preprocessing pipelines drive heterogeneous published methods to convergence on the released data, and (ii) the reported numbers can be reproduced bit-identically from the released Docker containers.

\todoTable{tab:baselines}{Test-split PSNR / SSIM / LPIPS, and per-task AUC for lung-nodule detection and Dice for liver-lesion segmentation, for four reproduced published baselines on PWM-LDCT. Each row is reproducible via the corresponding Docker container shipped with the release.}

\subsection*{Optional: framework-credential artifacts}
The release additionally ships 5-tuple credential JSON for each baseline under a companion signal-equivalence framework (manuscript in preparation,~\citealt{ws2framework}). These are optional reporting artifacts intended for credentialed leaderboard submission once the companion framework publishes; the descriptor's core validation evidence remains Table~\ref{tab:baselines}'s standard reconstruction-quality and task metrics, which do not depend on the framework. The credential set is summarized in Table~\ref{tab:credentials}; readers and reviewers who do not use the credential framework can skip this subsection.

\todoTable{tab:credentials}{Optional 5-tuple credential artifacts for four baselines at the canonical lung-nodule detection task, 25\% dose. Credential schema is defined in the companion framework manuscript (in preparation).}

%% =====================================================================
\section*{Usage Notes}

\subsection*{Loading the data}

Install the data loader:

\begin{verbatim}
    pip install pwm_ldct_loader
\end{verbatim}

Load a split:

\begin{verbatim}
    from pwm_ldct_loader import LowDoseCTDataset
    train = LowDoseCTDataset(root="/path/to/pwm_ldct",
                              split="train", seed=42)
    sample = train[0]
    # sample["full_dose"]   -> Tensor [H, W]
    # sample["low_dose"]    -> Tensor [H, W]  (real or sim)
    # sample["sinogram"]    -> Tensor [views, dets]
    # sample["annotations"] -> dict
    # sample["metadata"]    -> dict
\end{verbatim}

The loader pins random seeds, propagates them to PyTorch / NumPy / CUDA RNGs, and uses deterministic CuDNN settings.

\subsection*{Reproducing the baseline results}

\begin{verbatim}
    docker run --gpus all \
        -v /path/to/pwm_ldct:/data \
        -v $(pwd)/output:/output \
        pwm-ldct-red-cnn:v1 \
        eval --dataset /data --split test \
             --out /output/results.json --seed 42
\end{verbatim}

The output \texttt{results.json} should match the published values bit-identically (within IEEE~754 FP tolerance) on equivalent hardware.

\subsection*{Submitting to the permanent leaderboard}

Submissions to the PWM Low-Dose CT Challenge leaderboard (forthcoming; described in a companion manuscript in preparation,~\citealt{ws4leaderboard}) will take the form of a PWM RunBundle (Docker image + eval entry point + \texttt{results.json}, with optional 5-tuple credentials under the signal-equivalence framework in preparation,~\citealt{ws2framework}). The submission contract is documented at the leaderboard portal. Dataset users who do not wish to participate in the challenge can use the loader and Docker recipes directly --- the leaderboard is an ecosystem add-on, not a prerequisite for using the dataset.

\subsection*{Cross-vendor evaluation}

A method developer can request a held-out-vendor split via the loader's \texttt{leave\_vendor\_out=}\textit{vendor} argument; this filters the training split to exclude the named vendor and reassigns its patients to a separate evaluation slice. We recommend this protocol as the canonical cross-vendor evaluation; a companion reference-method manuscript exercises this protocol in detail (in preparation,~\citealt{ws3reference}).

\subsection*{Reuse cases}

We anticipate four primary reuse modes:

\begin{enumerate}
    \item \textbf{Reconstruction-method benchmarking.} Train on the released training split and evaluate on the released test split using the standard reconstruction-quality metrics provided by the loader (PSNR, SSIM, LPIPS, and per-task AUC / Dice). Methods that wish to participate in the companion benchmark may additionally submit credentials under the signal-equivalence framework and post to the permanent leaderboard once those ecosystem components launch (both in preparation,~\citealt{ws2framework,ws4leaderboard}).
    \item \textbf{Cross-vendor robustness studies.} Use the leave-one-vendor-out evaluation API to characterize how a method generalizes from a training-vendor distribution to a held-out-vendor distribution.
    \item \textbf{Simulation-vs-real noise modeling.} Use the paired real-low-dose and simulated-low-dose images on the \NvendorA{} patients with real paired acquisitions to characterize the discrepancy between common Poisson-thinning simulators and physical low-dose acquisition --- a long-standing gap in the field that this dataset's pairing structure now lets researchers quantify.
    \item \textbf{Downstream-task analysis.} The per-scan diagnostic-quality Likert labels and per-lesion annotations support studies that correlate computational reconstruction quality with radiologist-perceived diagnostic utility, beyond pixel-wise PSNR / SSIM.
\end{enumerate}

\subsection*{Analyses uniquely enabled by this dataset}

PWM-LDCT is not the first low-dose CT dataset; what makes it uniquely useful is the \emph{combination} of multi-vendor coverage with paired real-and-simulated low-dose at sufficient scale to support stratified statistical inference. Specific research questions that no existing public dataset can answer but that this one can:

\begin{enumerate}
    \item \textbf{How wrong is Poisson-thinning simulation, quantitatively, on real-clinical data, per vendor?} The community has long suspected that simulated low-dose underestimates real low-dose noise~\citep{wang2020deeplearningreview}, but per-vendor quantification at the $\geq 100$-patient scale has not been possible. The Mayo LDCT-PD release~\citep{mccollough2020mayo} permits this on Siemens hardware only; PWM-LDCT extends it to GE and Canon.
    \item \textbf{Does cross-vendor distribution shift mostly come from noise statistics or from reconstruction kernel?} The dataset's preserved per-vendor reconstruction kernels (with measured MTFs in metadata) plus the deliberate non-harmonization of noise enable ablation studies that disentangle these two sources.
    \item \textbf{Is uncertainty-quantification calibration vendor-dependent?} A reconstruction method trained on Siemens scans may produce well-calibrated UQ on Siemens but mis-calibrated UQ on GE. With paired ground truth across vendors and dose levels, this can be measured directly for the first time.
    \item \textbf{Does radiologist-perceived diagnostic confidence agree with computational reconstruction quality, and does that agreement depend on dose level?} The per-scan Likert labels plus the per-dose-level reconstructions permit a regression of human-perceived quality onto computational metrics with sufficient sample size to detect vendor-specific or anatomy-specific deviations.
    \item \textbf{What is the smallest dose ratio at which lung-nodule detection AUC remains within $\varepsilon$ of the full-dose reference, conditional on patient subpopulation?} PWM-LDCT provides the only public dataset with the multi-vendor paired-dose structure and the sample-size budget to answer this question with adequate statistical power; the companion signal-equivalence framework (in preparation,~\citealt{ws2framework}) is one possible formalism for reporting stratified evidence, though the analysis can also be conducted with standard paired-bootstrap confidence intervals on the released metrics.
\end{enumerate}

Each of these five questions is the basis for at least one follow-on publication that PWM-LDCT enables. Together they form the dataset's primary justification for the multi-vendor + paired-acquisition + content-addressed design choices, which individually would have been simpler but jointly are essential.

\subsection*{Limitations and bias acknowledgment}

A Data Descriptor that does not document the cohort's biases is incomplete. We acknowledge the following limitations of PWM-LDCT.

\paragraph{Geographic scope.} The cohort is recruited at \todo{N} clinical sites in \todo{regions}. While participating sites span academic and community settings, no European, Asian, or African sites are represented in v1.0.0. Generalization to international populations should be considered an open question, not a tested property. A v2 release will pursue international partnerships if v1 is well-received.

\paragraph{Vendor coverage.} The release covers $\geq$~2 of the four major CT manufacturers (Siemens, GE, Canon, Philips). Philips is currently underrepresented or absent (\todo{confirm at submission}). The leave-one-vendor-out evaluation API is technically functional regardless, but cross-vendor claims involving the absent vendor are extrapolation, not measurement, and should be reported as such.

\paragraph{Scanner generations.} All included scanners are within \todo{N} years of clinical deployment age at the time of acquisition. Methods evaluated against PWM-LDCT may degrade on older scanner generations (single-source, narrower detector arrays, earlier IR algorithms) that are still in clinical use globally but absent from this dataset.

\paragraph{Clinical indication.} The cohort is heavily weighted toward lung-cancer screening (chest) and oncology follow-up (abdomen). Other indications --- pulmonary embolism, trauma, pediatric --- are out of scope for v1. Methods optimized for PWM-LDCT may not transfer cleanly to those settings.

\paragraph{Pediatric population.} The cohort is adult-only ($\geq$~18 years). Pediatric CT, where dose reduction matters most~\citep{larson2018pediatric}, is the most-requested Year-2 extension and will be addressed in a separate companion dataset rather than a PWM-LDCT version increment.

\paragraph{Radiologist panel composition.} The annotation panel is drawn from \todo{N} institutions and includes \todo{N} fellowship-trained subspecialists. The panel does not include international radiologists, trainee radiologists, or non-subspecialist generalists; reader-study generalizations beyond fellowship-trained American academic radiologists are extrapolation.

\paragraph{Time of acquisition.} Scans were acquired between \todo{collection start} and \todo{collection end}. Subsequent changes to scanner software, IR algorithms, or clinical protocols at participating sites may make later real-world scans systematically different from the released cohort.

\paragraph{Selection bias in the prospective-consent subset.} The prospectively consented patients (those contributing real paired-dose acquisitions) self-selected by agreeing to a sequential second low-dose acquisition. We have no reason to suspect strong selection bias, but the consenting population is structurally different from a true random sample of CT patients at the participating sites. The retrospective-waiver cohort (full-dose only) is not subject to this selection effect.

These limitations are deliberate features of the v1 release rather than oversights; addressing each by extending the cohort would have delayed v1 indefinitely. We will track and respond to community feedback about which limitations matter most for the field and prioritize them in future versions per the policy in the Data Versioning section below.

%% =====================================================================
\section*{Data Versioning and Persistent Identifiers}

Each release of PWM-LDCT is uniquely identified by the SHA-256 hash of its \texttt{checksums.sha256} manifest, which itself hashes every released artifact. This content hash is the canonical scientific identifier of the release: any downstream user can locally verify that a downloaded copy matches the released hash, with no dependence on any single hosting institution remaining online.

\paragraph{Versioning scheme.} Releases follow a semantic \texttt{MAJOR.MINOR.PATCH} scheme:
\begin{itemize}
    \item \textbf{MAJOR} increments break dataset schema compatibility (e.g., metadata field renamed, annotation format changed). Downstream tooling pinned to a previous major version is not guaranteed to run.
    \item \textbf{MINOR} increments add new patients, new vendor coverage, new anatomy, or new annotation categories without breaking schema. Tooling pinned to a previous minor version continues to run; new content is available opt-in.
    \item \textbf{PATCH} increments correct errors (e.g., a scan that failed deferred QC, an annotation correction, a subject-withdrawal removal). Patch versions are bit-identical to the previous version except at the listed corrections.
\end{itemize}

\paragraph{Persistent-identifier anchors.} The release's content hash is mirrored to three independent anchors so that long-term citation resolves consistently:

\begin{enumerate}
    \item \textbf{PhysioNet DOI} (primary citation handle). Each release version is assigned a separate DOI on PhysioNet of the form \texttt{10.13026/\todo{prefix}-X.Y.Z}; downloads from that DOI return bit-identical content. This is the citation handle reviewers and readers are expected to use.
    \item \textbf{IPFS replica.} The PWM Protocol Foundation maintains an IPFS replica of the release content under the same SHA-256 content hash. This serves as a hosting-institution-independent retrieval path for the same content.
    \item \textbf{PWM L3 specification entry.} The release is also registered as an L3 specification (\texttt{pwm\_l3\_spec:pwm\_low\_dose\_ct\_benchmark:X.Y.Z}) on the PWM Protocol mainnet, which provides a third independent timestamp of the content hash for users who wish to audit version history programmatically.
\end{enumerate}

The scientific contribution is the content hash itself; the three anchors are interchangeable retrieval paths for the same hash, and the descriptor's reproducibility claims do not depend on any single anchor remaining available.

\paragraph{Errata, bug-reporting, and data-correction policy.} Errors discovered in a release (e.g., a missed PHI hit, a discordant annotation surviving QA, a Dockerfile reproducibility regression) are reported via GitHub issues at \url{https://github.com/integritynoble/low_dose_CT/issues} with the label \texttt{errata-v1.0}. Each report is triaged within \todo{N} business days; verified errors are batched into a PATCH release that increments the version, regenerates the content hash, and updates all three immutability anchors. Erratum notes (what changed, why, which scans are affected) are published as Supplementary Material alongside each PATCH release, and a cumulative errata changelog is maintained at \texttt{ERRATA.md} in the repository. Subscribers to the release mailing list are notified of every PATCH release. Prior PATCH releases remain resolvable via their content hashes so that papers citing an earlier version remain reproducible; we do not retract or rewrite historical versions.

\paragraph{Extension policy.} The release roadmap is driven by community feedback collected through PhysioNet's discussion forum and the leaderboard issue tracker. Standing extension priorities:

\begin{itemize}
    \item International sites (Year 2)
    \item Philips platform coverage (Year 2)
    \item Pediatric subpopulation (companion dataset, not a PWM-LDCT version increment)
    \item Older scanner generations (deferred unless community demand surfaces)
    \item Additional clinical indications --- PE, trauma, cardiac (Year 2--3)
\end{itemize}

A roadmap update is published at the start of each calendar year alongside the year's annual leaderboard review (in preparation,~\citealt{ws4leaderboard}), with specific commitments contingent on the year's funded extensions.

\paragraph{Long-term hosting commitment.} PhysioNet is committed to indefinite hosting of all credentialed-access datasets that pass its review; the IPFS pin under the same content hash provides a hosting-institution-independent retrieval path for derived non-PHI artifacts. The SHA-256 manifest hash --- resolvable through the PhysioNet DOI, the IPFS CID, or the PWM L3 registry entry --- serves as the canonical identifier of record regardless of which anchor is used.

%% =====================================================================
\section*{Code Availability}

\begin{itemize}
    \item Dataset loader: \url{https://pypi.org/project/pwm_ldct_loader/} (Apache 2.0)
    \item Preprocessing Dockerfiles: \url{https://github.com/integritynoble/low_dose_CT/tree/main/WS-1_dataset/pipelines/}
    \item Forward model used for low-dose simulation: \url{https://github.com/integritynoble/pwm/tree/main/packages/pwm_core/contrib/modalities/ct_radon.py}
    \item Manifest SHA-256 content hash: published in the PhysioNet release record. Mirrored anchors --- PhysioNet DOI \texttt{10.13026/\todo{prefix}-1.0.0}, IPFS CID \todo{CID}, and PWM L3 spec entry \texttt{pwm\_l3\_spec:pwm\_low\_dose\_ct\_benchmark:1.0.0} at \todo{contract address} --- all resolve to the same content.
\end{itemize}

%% =====================================================================
\section*{Author contributions}

\todo{Fill at submission. Use the CRediT taxonomy: Conceptualization, Methodology, Software, Validation, Formal analysis, Investigation, Resources, Data curation, Writing-Original Draft, Writing-Review \& Editing, Visualization, Supervision, Project administration, Funding acquisition.}

%% =====================================================================
\section*{Competing interests}

\todo{Declare at submission.}

%% =====================================================================
\section*{Acknowledgments}

The authors thank the radiologists who contributed annotations, the participating clinical sites, the PWM Protocol Foundation for infrastructure, and the open-source maintainers of \texttt{pwm\_core}, PyTorch, NumPy, and PyDICOM.

%% =====================================================================
\bibliographystyle{plainnat}
\bibliography{refs}

\end{document}
